\section{Related Work}
\label{sec:related}

The setting we study lies at the intersection of three lines of work that have largely developed in
isolation: how federated learning communicates, how it is made private, and how it has been compressed into
a single round. Tracing how each arrived at its present state makes clear both what is established and where
the gap we address lies.

Federated learning began as an iterative procedure. In the original formulation a server repeatedly
broadcasts a global model, clients improve it on their local data, and the server averages the returned
updates~\cite{mcmahan2017communication}. This works well but inherits two persistent costs from its
iterative nature. The first is communication: reaching a good model can take tens or hundreds of rounds,
each moving a full model in both directions. The second is exposure: every round reveals a fresh update
computed on private data, and a long sequence of such updates is a rich target for inference. Much of the
subsequent literature can be read as attempts to control one of these two costs without worsening the
other.

The exposure problem has been met with two broad tools, distinguished by whether they alter the model that
is ultimately released. Secure aggregation lets a server learn only the sum of the clients' updates and
nothing about any individual one, using masking protocols that cancel in the aggregate~\cite
{bonawitz2017practical}; it is lossless, in that the aggregate is exact, but it protects only the per-round
sum and still operates within the iterative protocol. Differential privacy instead perturbs the updates
themselves, most commonly through differentially private stochastic gradient descent~\cite{abadi2016deep};
the guarantee is strong and composable, but because the noise enters the released model it trades privacy
directly against accuracy. This distinction, between privacy that costs accuracy and privacy that does not, is the axis along which our own privacy
design is best understood, and it recurs throughout what follows.

Homomorphic encryption offers a third route, protecting the updates by computing on them while they remain
encrypted. In its most ambitious form it has been used to train a neural network entirely under encryption
across many parties. POSEIDON, for instance, runs encrypted stochastic gradient descent under multiparty
CKKS, with the nonlinear activations replaced by polynomial approximations so that they can be evaluated
homomorphically~\cite{sav2021poseidon}. The privacy adds no accuracy cost and the result is strong, but the approach
remains fundamentally iterative and pays for the encrypted nonlinearities with a deep multiplicative
circuit, repeated bootstrapping, and a wall-clock measured in hours. A lighter line keeps the encryption
but narrows its job to aggregation alone: schemes that combine multi-key CKKS with masking reduce a round
of secure summation to a single non-interactive client-to-server transmission~\cite{kemmaka2025hybagg}.
These are efficient and genuinely cryptographic, but they encrypt the sum of model updates within an
iterative training protocol rather than a one-shot contribution, and they perform aggregation rather than
knowledge transfer. On the encryption axis these works are the closest to ours, and the contrast is
precise: they are either multi-round or are secure-aggregation primitives, whereas we encrypt a single
fine-tuning displacement.

Homomorphic encryption has also been applied to federated learning purely for aggregation, encrypting the
per-round updates so the server can sum them without seeing any in the clear; BatchCrypt, FedML-HE, and
FedSHE develop this for cross-silo training with various packing and quantization
optimisations~\cite{zhang2020batchcrypt,jin2023fedml,wei2025fedshe}, and a recent line encrypts only the
low-rank adapters of a parameter-efficient fine-tune: SHE-LoRA selectively encrypts a sensitivity-chosen
subset of the adapter under CKKS each round~\cite{li2025shelora}. These remain multi-round and, like
secure aggregation, protect the sum rather than transferring knowledge; encrypting the adapter rather than
the full model narrows the per-round object, but the cost is still paid on every round of an iterative
protocol, and the bilinear coupling of the two adapter factors (discussed below) is left to be resolved
under encryption. Our protocol differs on both axes: a single round, and an adapter parameterisation for
which the aggregate is exactly linear, so no coupling remains to resolve under encryption. A recurring obstacle across all
encrypted-training work is the nonlinearity: activations must be replaced by polynomial approximations to be
evaluated homomorphically, which deepens the multiplicative circuit and is delicate to train, as a line of
work on polynomial activations and encryption-friendly networks documents~\cite{garimella2021sisyphus,baruch2022methodology,agamennone2025polynomial,alhossain2025training,ibarrondo2021fhebn}.
Our design sidesteps this entirely: the only encrypted operation is linear, so no activation is ever
evaluated under encryption and the circuit stays at depth one.

Closest to our object, \emph{transfer learning} itself has been placed under encryption, though in settings
structurally different from ours. HETAL trains a softmax classifier head on a frozen public backbone entirely
under CKKS, presented as the first practical scheme for encrypted training of such a head~\cite{lee2023hetal};
a longer line trains linear and logistic models under homomorphic encryption by iterative encrypted gradient
descent~\cite{kim2018securelr,kim2018logistic}, recently extended to small networks~\cite{chiang2025cnn}. All
of these run an \emph{optimiser on ciphertexts}: the multiplicative depth grows with the number of steps, the
softmax and sigmoid nonlinearities must be polynomial-approximated, and bootstrapping is often required; and
they are single-party --- one data owner outsourcing computation --- rather than a collaboration among
mutually distrustful parties. A federated variant, Priv-FedTL, encrypts only the fine-tuned last layer of a
frozen backbone across clients, but aggregates it over many rounds of federated
training~\cite{privfedtl2026}; a concurrent scheme homomorphically averages encrypted feature tokens for
encrypted inference under a single decryption key~\cite{alamin2025vit}. We differ from all of these on the
axis that sets the cost: we perform \emph{no} encrypted training. Each client fine-tunes in plaintext on its
own data, and the server's only homomorphic operation is a single depth-one weighted sum of the encrypted
displacements, decrypted by a threshold of clients --- no encrypted optimiser, no polynomial nonlinearity, no
bootstrapping, and no single party trusted to train or to decrypt.

A separate line federates \emph{parameter-efficient fine-tuning}, communicating only a low-rank adapter
rather than the full model. FedIT first applied this to instruction tuning by averaging the adapters round
by round~\cite{zhang2024fedit}, but naive per-factor averaging is biased: a client learns an update
$\Delta W = BA$, and averaging the factors separately yields $(\sum_j B_j)(\sum_j A_j)\neq\sum_j B_jA_j$,
a discrepancy formalised as aggregation noise by FLoRA, which removes it by stacking the clients'
factors~\cite{wang2024flora}; FlexLoRA instead reconstructs and averages the full products before
redistributing by SVD~\cite{bai2024flexlora}, HetLoRA reconciles heterogeneous ranks~\cite{cho2024hetlora},
and FedSA-LoRA shares only one factor~\cite{guo2025fedsalora}. The fix we adopt is FFA-LoRA, which freezes
the randomly initialised down-projection and trains only the up-projection, so that averaging the trained
factor is exact; it was introduced to stabilise multi-round, differentially private federated
tuning~\cite{sun2024ffalora}. Under encryption this choice has a sharper consequence than in plaintext: the
bilinearity is not merely a bias to correct but a cost cliff, since resolving it on the server would require
ciphertext-by-ciphertext products, whereas freezing the down-projection makes the one-shot encrypted
aggregate exact at multiplicative depth one. Our aggregation is, algebraically, task-vector
merging~\cite{ilharco2023editing}, which has been shown equivalent to one-shot federated
averaging~\cite{tao2024taskarith}; that a single round of fine-tuning suffices for foundation models is
itself established in plaintext~\cite{malinovsky2024oneround}. What none of these provide is cryptographic
protection of the contributions, which is our subject.

The communication problem has driven the parallel development of one-shot federated learning, which
collapses the entire exchange into a single upload and download~\cite{guha2019oneshot,wang2025towards}. The earliest approaches transferred
knowledge rather than weights, having clients exchange predictions on a shared public set so that
heterogeneous models could teach one another~\cite{li2019fedmd,itahara2021dsfl}; in practice these
remained iterative, repeating the exchange to converge. Ensemble distillation made the distillation step
explicit by training a server model against the clients' aggregated predictions, though in its standard
configuration it too runs for many rounds~\cite{lin2020feddf}. The decisive step to a true single round
came with data-free one-shot methods, which avoid any shared data by having the server distil the clients'
models through a generator or through synthesised inputs, as in DENSE and its
successors~\cite{zhang2022dense,dai2024coboosting}, and with fusion-based methods that stitch client models
together with lightweight learned adapters~\cite{tang2024fusefl} or aggregate them through a layer-wise
posterior~\cite{liu2024fedlpa,zhang2024fedsd2c}. These methods are the closest to ours on the one-shot and
distillation axes, and they establish that a single round can produce a usable joint model. What they do
not provide is any cryptographic protection: they operate in plaintext, so the contributions each client
sends are exposed.

Closing that last gap, a handful of works add privacy to one-shot federated learning, and they do so with
differential privacy. FedAUXfdp distils through a pretrained feature extractor on auxiliary public data and
releases the client contributions under a differentially private mechanism~\cite{hoech2022fedauxfdp};
related approaches protect a diffusion-generated~\cite{feddiff2024} or a teacher-ensemble~\cite
{li2021fedkt} surrogate, or add noise-free differential privacy to the distillation
itself~\cite{sun2021fedmdnfdp}, under the same kind of guarantee. These are the natural quantitative peers
for our method, since they share its one-shot, distillation-based structure, but their privacy is lossy in
exactly the sense above: the noise needed for a meaningful budget is added to the artifact that is released,
and accuracy falls as the budget tightens.

This lossy trade-off is most sharply visible in differentially private transfer learning, the lossy
counterpart of our own setting. A well-established line shows that differential privacy works best precisely
on \emph{frozen} public features: linear or parameter-efficient models trained with DP-SGD on a pretrained
backbone reach strong accuracy at moderate privacy
budgets~\cite{tramer2021features,mehta2022dpfeatures,yu2022dpft,li2022dplearners}, and this extends to
federated low-rank adaptation under differential privacy~\cite{liu2023dplora,xu2024dpdylora}. These methods
share our backbone-plus-small-unit structure but pay a budget-dependent accuracy tax, since the noise enters
the released model, and they are either centralized or run over many federated rounds. Our protocol occupies
the same structural niche without the tax: the contributions are protected cryptographically rather than by
perturbation, the released model carries no privacy-induced noise, and the entire exchange is a single round.
To our knowledge no prior method is simultaneously one-shot, federated, and private --- whether by
homomorphic encryption or by differential privacy --- in this transfer-learning setting.

The reason privacy is the binding concern here is that the artifacts federated protocols exchange leak.
Sharing gradients or model updates permits reconstruction of training inputs~\cite{zhu2019deep,dimitrov2024spear,carletti2025sok}
and supports membership and property inference~\cite{shokri2017membership,nasr2019comprehensive,yeom2018privacy,salem2019ml,carlini2022membership,kerkouche2023property},
and even released confidences enable model inversion~\cite{fredrikson2015model,hitaj2017deep}. Federated
distillation, which shares predictions or a public-data signal instead of gradients, was hoped to sidestep
this, but a growing body of attacks shows it leaks as well: paired-logit inversion recovers images from
shared logits~\cite{takahashi2023breaching}, and membership can be inferred from distillation outputs and
from public-dataset usage~\cite{yang2023fdleaks,wang2024graddiff,jagielski2023students,gu2023ldia,shi2025unveiling}.
This is why we protect the contributions cryptographically rather than by perturbation, and why we measure
in \cref{sec:mia} what the released model still reveals.

Seen together, the two strongest lines never meet. The cryptographic work that would protect contributions
without loss is either multi-round encrypted training or a secure-aggregation primitive, and the one-shot
work that communicates in a single round is either plaintext or made private only by lossy
perturbation. The method we present occupies the empty intersection: a one-shot federated learning protocol
whose single encrypted operation is a linear aggregation under multiparty homomorphic encryption, giving
the accuracy-preserving privacy of the cryptographic line at the communication budget of the one-shot line.
